\documentclass[exercice]{thedocument}%
\startdatakeys
\source{}
\cycle{}
\competencesDuSocle{}
\connaissancesRequises{}
\competencesTravaillees{}
\enddatakeys
\begin{document}%
\!S3coord=0,1;%
\!S1coord=-3,0;%
\!S5coord=3,0;%
\!S6coord=0,-3;%
\begin{flushleft}
    Calculer le périmètre du polygone \<S1;\<S2;\<S3;\<S4;\<S5;\<S6;.\par
    On donnera un arrondi au dixième du résultat.
\end{flushleft}\hspace{0.03\linewidth}
\begin{center}
    \begin{tikzpicture}[scale=0.7]
        \draw[line width=0.5 pt](-4,-4)grid(4,4);
        \coordinate[label=above:\<S3;](S3)at(\<S3coord;);
        \coordinate[label=below:\<S6;](S6)at(\<S6coord;);
        \coordinate[label=left:\<S1;](S1)at(\<S1coord;);
        \coordinate[label=\<S2;](S2)at(-2,2);
        \coordinate[label=above:\<S4;](S4)at(2,2);
        \coordinate[label=right:\<S5;](S5)at(\<S5coord;);
        \draw(S1)--(S2)--(S3)--(S4)--(S5)--(S6)--cycle;
    \end{tikzpicture}
\end{center}
\end{document}
